\includegraphics[width=3.5in,height=1.15625in]{media/image1.png}

\textbf{School of Computer Science Engineering and Information Systems
(SCORE)}

\textbf{Winter Semester 2024-2025}

\textbf{ISWE207P-OBJECT ORIENTED ANALYSIS AND DESIGN LAB}

System Design Document

Bank Management System

\textbf{HARSHITH B 23MIS0012}

\textbf{L45 + L46}

\textbf{Associate Professor Sr Grade 2 Dr Magesh}

\begin{longtable}[]{@{}
  >{\raggedright\arraybackslash}p{(\linewidth - 2\tabcolsep) * \real{0.4364}}
  >{\raggedright\arraybackslash}p{(\linewidth - 2\tabcolsep) * \real{0.5636}}@{}}
\toprule\noalign{}
\endhead
\bottomrule\noalign{}
\endlastfoot
Document ID & SYSTEM DESIGN-v0.1 \\
Version Number & \\
Issue Date & December 13, 2024 \\
Classification & Public \\
\end{longtable}

Copyright Notice

©Trojan Solutions Limited, (2024 -- 2025)

All Rights Reserved

The information contained in this document is the property of Trojan
Solutions Limited. No part of this document may be reproduced, stored in
a retrieval system, or transmitted in any form, or by any means;
mechanical, photocopying, recording, or otherwise, without the prior
written consent of Trojan Solutions Limited. Under the law, copying
includes translating into another language or format. Legal action will
be taken against any infringement.

The information contained in this document is subject to change without
notice and does not carry any contractual obligation for Trojan
Solutions Limited. Trojan Solutions Limited reserves the right to make
changes to any products or services described in this document at any
time without notice. Trojan Solutions Limited shall not be held
responsible for the direct or indirect consequences of the use of the
information contained in this document.

\textbf{Revision History}

\begin{longtable}[]{@{}
  >{\raggedright\arraybackslash}p{(\linewidth - 6\tabcolsep) * \real{0.2263}}
  >{\raggedright\arraybackslash}p{(\linewidth - 6\tabcolsep) * \real{0.2070}}
  >{\raggedright\arraybackslash}p{(\linewidth - 6\tabcolsep) * \real{0.3098}}
  >{\raggedright\arraybackslash}p{(\linewidth - 6\tabcolsep) * \real{0.2568}}@{}}
\toprule\noalign{}
\endhead
\bottomrule\noalign{}
\endlastfoot
\textbf{Date} & \textbf{Version} & \textbf{Description} & \textbf{Author
(s)} \\
/01/2020 & & Draft Version & Harshith B \\
& & & \\
& & & \\
\end{longtable}

\begin{longtable}[]{@{}
  >{\raggedright\arraybackslash}p{(\linewidth - 4\tabcolsep) * \real{0.5892}}
  >{\raggedright\arraybackslash}p{(\linewidth - 4\tabcolsep) * \real{0.2595}}
  >{\raggedright\arraybackslash}p{(\linewidth - 4\tabcolsep) * \real{0.1512}}@{}}
\toprule\noalign{}
\endhead
\bottomrule\noalign{}
\endlastfoot
\textbf{Reviewed By (Customer)} & \textbf{Signature} & \textbf{Date} \\
& & \\
& & \\
& & \\
& & \\
& & \\
\end{longtable}

The reviewer signoff shall signify the recommendation for acceptance of
this document.

\begin{longtable}[]{@{}
  >{\raggedright\arraybackslash}p{(\linewidth - 2\tabcolsep) * \real{0.5403}}
  >{\raggedright\arraybackslash}p{(\linewidth - 2\tabcolsep) * \real{0.4597}}@{}}
\toprule\noalign{}
\begin{minipage}[b]{\linewidth}\raggedright
\textbf{Prepared By}
\end{minipage} & \begin{minipage}[b]{\linewidth}\raggedright
\textbf{Acknowledged By}
\end{minipage} \\
\midrule\noalign{}
\endhead
\bottomrule\noalign{}
\endlastfoot
& \\
Mr. Harshith B & Mr. Kaarthik. M \\
Title: Chief Requirements Officer & Title: Chief Software Officer \\
Trojan Solutions Limited & Vellore Limited \\
Date: 25/03/2025 & Date: 25/03/2025 \\
\end{longtable}

\begin{longtable}[]{@{}
  >{\raggedright\arraybackslash}p{(\linewidth - 2\tabcolsep) * \real{0.5015}}
  >{\raggedright\arraybackslash}p{(\linewidth - 2\tabcolsep) * \real{0.4985}}@{}}
\toprule\noalign{}
\begin{minipage}[b]{\linewidth}\raggedright
\textbf{Accepted By}
\end{minipage} & \begin{minipage}[b]{\linewidth}\raggedright
\textbf{Accepted By}
\end{minipage} \\
\midrule\noalign{}
\endhead
\bottomrule\noalign{}
\endlastfoot
& \\
Mr. Nandakrishnan & Mr. Rithik. G \\
Title: Chied Technology Officer & Title: Senior Architect \\
Google Tech Solutions Limited & Google Tech Solutions Limited \\
Date: 26/03/2025 & Date: 26/03/2025 \\
\end{longtable}

\textbf{Table of Contents}

\hyperref[list-of-tables]{List of Tables 5}

\hyperref[list-of-figures]{List of Figures
\hyperref[list-of-figures]{6}}

\hyperref[introduction]{1 Introduction \hyperref[introduction]{7}}

\hyperref[purpose-of-document]{1.1 Purpose of Document
\hyperref[purpose-of-document]{7}}

\hyperref[X918f79d1401d26e11fe5c5060cfda901dfd4841]{1.2 Document Scope
\hyperref[X918f79d1401d26e11fe5c5060cfda901dfd4841]{8}}

\hyperref[in-scope]{1.2.1 In-Scope \hyperref[in-scope]{8}}

\hyperref[out-of-scope]{1.2.2 Out-of-Scope \hyperref[out-of-scope]{8}}

\hyperref[assumptions]{1.2.3 Assumptions \hyperref[assumptions]{8}}

\hyperref[methodology-tools-and-approach]{1.3 Methodology, Tools, and
Approach \hyperref[methodology-tools-and-approach]{8}}

\hyperref[acronyms-and-abbreviations]{1.4 Acronyms and Abbreviations
\hyperref[acronyms-and-abbreviations]{8}}

\hyperref[design-overview]{2 Design Overview
\hyperref[design-overview]{9}}

\hyperref[background-information]{2.1 Background Information
\hyperref[background-information]{9}}

\hyperref[system-evolution-description]{2.2 System Evolution Description
\hyperref[system-evolution-description]{9}}

\hyperref[required-environment]{2.3 Required Environment
\hyperref[required-environment]{10}}

\hyperref[Xb309f13310ad76dd529681052fe5315a8c7c065]{2.4 Constraints
\hyperref[Xb309f13310ad76dd529681052fe5315a8c7c065]{\textbf{.}}}10

\hyperref[Xfce2f82cfd63ce3cc56ff7cd87639809d1de288]{2.5 Design
Trade-offs \hyperref[Xfce2f82cfd63ce3cc56ff7cd87639809d1de288]{11}}

\hyperref[structural-family-diagrams]{3 Structural Family Diagrams
\hyperref[structural-family-diagrams]{11}}

\hyperref[X064e3884438d5a452559189fc7e45d065967ef0]{3.1 Class Diagram
11}

\hyperref[X74cee0fdbef85b50ea600fc9189276a3bfcd0e2]{3.2 Object Diagram
14}

\hyperref[X1697ccb2e76750f21357dbfadeb6e81b5e54e78]{3.3 Package Diagram
15}

\hyperref[behavioral-family-diagrams]{4 Behavioral Family Diagrams 15}

4.1 Use Case Diagram 15

4.2 Use Case Descriptions 15

4.3 State Chart Diagram 16

4.4 Activity Diagram 16

4.5 Sequence Diagram 16

4.6 Collaboration Diagram 17

\hyperref[X2415f40c9d0f9e87bdbc628763cbc47606accb8]{5 Architectural
Family Diagrams 17}

\hyperref[X131ff30b2dee961d64957eddfef9aef7e4844bb]{5.1 Component
Diagram 17}

\hyperref[X40e4a2f891c91a91e35c42386e697b8abc32e12]{5.2 Deployment
Diagram 17}

\hyperref[appendix]{Appendix 18}

\section{List of Tables}\label{list-of-tables}

\begin{enumerate}
\def\labelenumi{\arabic{enumi}.}
\item
  Acronyms and Abbreviations
\item
  Requirements Traceability Matrix
\end{enumerate}

\section{List of Figures}\label{list-of-figures}

\begin{enumerate}
\def\labelenumi{\arabic{enumi}.}
\item
  Class Diagram
\end{enumerate}

\section{Introduction}\label{introduction}

The \textbf{Bank Management System} is a software application designed
to manage and automate core banking operations efficiently. The system
provides a centralized platform for handling various banking services
such as customer account management, transaction processing, loan
handling, and secure authentication for customers and employees. By
integrating digital solutions, the system minimizes human errors,
enhances operational efficiency, and ensures seamless financial
transactions.

With the increasing reliance on digital banking services, a
well-structured Bank Management System plays a crucial role in improving
service accessibility, security, and compliance with banking
regulations. The system is designed to support both online and offline
banking operations, ensuring smooth interaction between customers and
financial institutions. It offers functionalities such as automated
transaction processing, real-time balance updates, detailed reporting,
and enhanced security protocols to protect sensitive financial data.

By incorporating modern technologies such as cloud computing, data
encryption, and role-based access control, the system enhances banking
services while maintaining industry standards. This document provides an
in-depth view of the system's design, methodologies, and implementation
strategies to ensure a robust and scalable banking solution.

\subsection{Purpose of Document}\label{purpose-of-document}

\protect\phantomsection\label{X918f79d1401d26e11fe5c5060cfda901dfd4841}{}The
purpose of this document is to provide a structured and detailed
overview of the \textbf{Bank Management System} and its design. It
outlines the systems architecture, key functionalities, and
implementation strategies, ensuring clarity in development and
deployment.

This document serves as a foundational guide for developers, project
managers, system architects, and stakeholders to understand how the
system is structured and how it aligns with the requirements of modern
banking operations. By providing insights into the software and hardware
components, database structure, and user interface design, it
facilitates seamless integration and future scalability.

Additionally, it ensures that security, efficiency, and compliance with
financial regulations are well integrated into the system's design. The
document also defines various performance benchmarks, best practices,
and industry standards that must be adhered to throughout the system's
lifecycle. Through comprehensive documentation, this SDD aims to create
a structured roadmap for successful implementation, maintenance, and
enhancement of the \textbf{Bank Management System} over time.

\subsection{Document Scope}\label{document-scope}

This document covers the systems design, including structural and
behavioral diagrams, constraints, and trade-offs involved in the
implementation.

\subsubsection{In-Scope}\label{in-scope}

\begin{itemize}
\item
  Customer Account Management
\item
  Transaction Handling (Deposits, Withdrawals, Transfers)
\item
  Loan Management
\item
  User Authentication and Role-Based Access
\item
  Report Generation
\end{itemize}

\subsubsection{Out-of-Scope}\label{out-of-scope}

\begin{itemize}
\item
  Third-party Payment Gateway Integrations
\item
  AI-based fraud detection (Future Enhancement)
\item
  Advanced Financial Forecasting
\end{itemize}

\subsubsection{Assumptions}\label{assumptions}

\begin{itemize}
\item
  The system will be used by bank employees and customers.
\item
  Internet connectivity is available for online banking features.
\item
  Security protocols such as encryption and authentication will be in
  place.
\end{itemize}

\subsection{Methodology, Tools, and
Approach}\label{methodology-tools-and-approach}

\begin{itemize}
\item
  \textbf{\ul{Methodology}:} Object-Oriented Analysis and Design (OOAD)
\item
  \textbf{\ul{Tools Used}:} UML Diagrams (Lucidchart), MySQL, Spring
  Boot, Java, React.js
\item
  \textbf{\ul{Approach}:} Agile development methodology
\end{itemize}

\subsection{Acronyms and
Abbreviations:}\label{acronyms-and-abbreviations}

\begin{longtable}[]{@{}
  >{\centering\arraybackslash}p{(\linewidth - 2\tabcolsep) * \real{0.1824}}
  >{\centering\arraybackslash}p{(\linewidth - 2\tabcolsep) * \real{0.8176}}@{}}
\toprule\noalign{}
\begin{minipage}[b]{\linewidth}\centering
GUI
\end{minipage} & \begin{minipage}[b]{\linewidth}\centering
Graphical User Interface
\end{minipage} \\
\midrule\noalign{}
\endhead
\bottomrule\noalign{}
\endlastfoot
SDD & System Design Document \\
API & Application Programming Interface \\
DBMS & Database Management System \\
RBAC & Role-Based Access Control \\
TLS & Transport Layer Security \\
KYC & Know Your Customer \\
AML & Anti-Money Laundering \\
\end{longtable}

\section{Design Overview}\label{design-overview}

\subsection{Background Information}\label{background-information}

The \textbf{Bank Management System} is developed to address the growing
need for an efficient, secure, and scalable platform for handling
banking operations. Traditional banking systems rely heavily on manual
processes, leading to inefficiencies, delays, and security
vulnerabilities. The proposed system automates and streamlines key
banking functions such as account management, transaction processing,
and loan management, ensuring a seamless experience for both customers
and banking staff.

The system is designed to operate in an online environment, providing
customers with 24/7 access to banking services. It integrates with
financial institutions databases and APIs to ensure real-time data
processing, secure transactions, and compliance with banking
regulations. The Bank Management System aims to enhance customer
satisfaction by offering a user-friendly interface, quick transaction
processing, and robust security measures.

\subsection{System Evolution
Description}\label{system-evolution-description}

Banking systems have evolved from traditional ledger-based accounting to
digital solutions with automated workflows. The \textbf{Bank Management
System} builds upon existing digital banking frameworks by incorporating
advanced security features, cloud-based storage, and enhanced user
authentication mechanisms. The system also allows for seamless
integration with third-party financial services such as payment
gateways, mobile banking applications, and regulatory compliance
platforms.

Improvements over traditional banking systems include:

\begin{itemize}
\item
  Real-time transaction updates to reduce processing delays.
\item
  Secure authentication protocols (e.g., multi-factor authentication and
  biometric verification).
\item
  Enhanced reporting and analytics for improved financial
  decision-making.
\item
  Integration with mobile banking and digital payment solutions.
\end{itemize}

\subsection{Required Environment}\label{required-environment}

To ensure the smooth operation of the \textbf{Bank Management System},
the following environments are required:

\begin{longtable}[]{@{}
  >{\raggedright\arraybackslash}p{(\linewidth - 2\tabcolsep) * \real{0.2829}}
  >{\raggedright\arraybackslash}p{(\linewidth - 2\tabcolsep) * \real{0.7171}}@{}}
\toprule\noalign{}
\endhead
\bottomrule\noalign{}
\endlastfoot
Product/Solution & Environment \\
Bank Management System & - \textbf{Production} (Live environment for
real-time banking operations) \\
& - \textbf{Development} (Used for testing and feature development) \\
& - \textbf{Staging} (Pre-production environment for final testing
before deployment) \\
& - \textbf{Backup \& Recovery} (Dedicated environment for data backups
and disaster recovery) \\
\end{longtable}

\subsection{This structured environment ensures the system remains
scalable, secure, and efficient while supporting continuous improvements
and compliance with banking
regulations.}\label{this-structured-environment-ensures-the-system-remains-scalable-secure-and-efficient-while-supporting-continuous-improvements-and-compliance-with-banking-regulations.}

\subsection{Constraints}\label{constraints}

\protect\phantomsection\label{Xfce2f82cfd63ce3cc56ff7cd87639809d1de288}{}The
design of the \textbf{Bank Management System} is influenced by various
constraints that impact on its functionality, performance, and
scalability:

\begin{itemize}
\item
  \textbf{Hardware and Software Environment:} The system must run on
  secure, high-performance servers with support for modern web
  technologies.
\item
  \textbf{End-User Environment:} The system must be accessible across
  various devices, including desktops, tablets, and mobile phones.
\item
  \textbf{Resources Availability:} Database and API services must be
  consistently available to prevent downtime.
\item
  \textbf{Standards Compliance:} The system must adhere to banking and
  financial regulatory standards such as PCI DSS, GDPR, and RBI
  guidelines.
\item
  \textbf{Interoperability Requirements:} The system should integrate
  seamlessly with third-party banking APIs and payment gateways.
\item
  \textbf{Security Requirements:} Data encryption, role-based access
  control, and fraud detection mechanisms must be implemented.
\item
  \textbf{Performance Requirements:} The system should be capable of
  handling a large volume of transactions without performance
  degradation.
\item
  \textbf{Network Communications:} Secure and efficient communication
  protocols should be used to protect data transmission over the
  network.
\item
  \textbf{Verification and Validation Requirements:} Rigorous testing,
  including unit testing, integration testing, and security testing,
  must be conducted before deployment.
\end{itemize}

\subsection{Design Trade-offs}\label{design-trade-offs}

Design decisions often involve trade-offs to balance performance,
security, and usability. The \textbf{Bank Management System} considers
the following trade-offs:

\begin{itemize}
\item
  \textbf{Security vs. User Experience:} Implementing strong security
  measures (e.g., multi-factor authentication) may slightly increase
  login time but ensures data protection.
\item
  \textbf{Performance vs. Cost:} Cloud-based deployment offers
  scalability but incurs higher operational costs compared to
  on-premises hosting.
\item
  \textbf{Feature Set vs. Simplicity:} Including advanced features such
  as AI-driven fraud detection enhances security but increases system
  complexity.
\item
  \textbf{Real-time Processing vs. System Load:} Enabling real-time
  transaction updates improves user experience but requires more
  computational resources.
\end{itemize}

By carefully evaluating these trade-offs, the system ensures a balance
between efficiency, security, and ease of use, providing a robust
banking solution that meets modern financial requirements.

\section{\texorpdfstring{Structural Family Diagrams
}{Structural Family Diagrams }}\label{structural-family-diagrams}

\subsection{Class Diagram}\label{class-diagram}

\section{\texorpdfstring{ Classes}{ Classes}}\label{classes}

Each entity in the banking system is modeled as a \textbf{class} with
attributes and methods defining its behavior.

\subsection{Customer}\label{customer}

The \texttt{Customer} class represents bank customers and includes the
following attributes and methods:

\begin{itemize}
\item
  \textbf{Attributes:}

  \begin{itemize}
  \item
    \texttt{customerID}: Unique identifier for the customer.
  \item
    \texttt{name}: Full name of the customer.
  \item
    \texttt{address}: Residential address of the customer.
  \item
    \texttt{accountList}: List of accounts associated with the customer.
  \end{itemize}
\item
  \textbf{Methods:}

  \begin{itemize}
  \item
    \texttt{register(}\texttt{)} - Registers a new customer.
  \item
    \texttt{updateProfile}\texttt{(}\texttt{)} - Updates customer
    information.
  \end{itemize}
\end{itemize}

\subsection{BankEmployee}\label{bankemployee}

The \texttt{BankEmployee} class represents employees managing banking
operations.

\begin{itemize}
\item
  \textbf{Attributes:}

  \begin{itemize}
  \item
    \texttt{employeeID}: Unique identifier for the employee.
  \item
    \texttt{name}: Name of the employee.
  \item
    \texttt{role}: Job role (e.g., teller, manager).
  \end{itemize}
\item
  \textbf{Methods:}

  \begin{itemize}
  \item
    \texttt{processTransaction}\texttt{(}\texttt{)} - Authorizes
    transactions.
  \item
    \texttt{approveLoan}\texttt{(}\texttt{)} - Approves or rejects loan
    applications.
  \end{itemize}
\end{itemize}

\subsection{Account}\label{account}

The \texttt{Account} class serves as a base class with specializations
such as \texttt{SavingsAccount} and \texttt{CurrentAccount}.

\begin{itemize}
\item
  \textbf{Attributes:}

  \begin{itemize}
  \item
    \texttt{accountNumber}: Unique identifier for the account.
  \item
    \texttt{balance}: Current balance in the account.
  \item
    \texttt{accountType}: Specifies account type (Savings/Current).
  \end{itemize}
\item
  \textbf{Methods:}

  \begin{itemize}
  \item
    \texttt{deposit(amount)} - Deposits money into the account.
  \item
    \texttt{withdraw(amount)} - Withdraws money from the account.
  \item
    \texttt{checkBalance}\texttt{(}\texttt{)} - Returns the current
    balance.
  \end{itemize}
\end{itemize}

\subsection{Transaction}\label{transaction}

The \texttt{Transaction} class represents any financial operation such
as deposit, withdrawal, and transfer.

\begin{itemize}
\item
  \textbf{Attributes:}

  \begin{itemize}
  \item
    \texttt{transactionID}: Unique identifier for the transaction.
  \item
    \texttt{transactionType}: Type of transaction
    (Deposit/Withdrawal/Transfer).
  \item
    \texttt{amount}: Amount involved in the transaction.
  \item
    \texttt{timestamp}: Date and time of the transaction.
  \end{itemize}
\item
  \textbf{Methods:}

  \begin{itemize}
  \item
    \texttt{processTransaction}\texttt{(}\texttt{)} - Executes the
    transaction.
  \item
    \texttt{validateTransaction}\texttt{(}\texttt{)} - Validates the
    transaction details.
  \end{itemize}
\end{itemize}

\subsection{Loan}\label{loan-1}

The \texttt{Loan} class represents loan-related functionalities such as
personal, home, or car loans.

\begin{itemize}
\item
  \textbf{Attributes:}

  \begin{itemize}
  \item
    \texttt{loanID}: Unique identifier for the loan.
  \item
    \texttt{loanType}: Type of loan (Home/Car/Personal).
  \item
    \texttt{loanAmount}: The borrowed amount.
  \item
    \texttt{interestRate}: Interest rate applied to the loan.
  \item
    \texttt{tenure}: Loan repayment period.
  \end{itemize}
\item
  \textbf{Methods:}

  \begin{itemize}
  \item
    \texttt{approveLoan}\texttt{(}\texttt{)} - Approves or rejects a
    loan.
  \item
    \texttt{calculateEMI}\texttt{(}\texttt{)} - Calculates the monthly
    installment amount.
  \end{itemize}
\end{itemize}

\subsection{Payment}\label{payment}

The \texttt{Payment} class handles payments for loans or bills.

\begin{itemize}
\item
  \textbf{Attributes:}

  \begin{itemize}
  \item
    \texttt{paymentID}: Unique identifier for the payment.
  \item
    \texttt{amount}: The amount paid.
  \item
    \texttt{paymentMethod}: Mode of payment (Cash/Card/Online).
  \item
    \texttt{status}: Payment status (Success/Failed/Pending).
  \end{itemize}
\item
  \textbf{Methods:}

  \begin{itemize}
  \item
    \texttt{processPayment}\texttt{(}\texttt{)} - Executes the payment.
  \item
    \texttt{refundPayment}\texttt{(}\texttt{)} - Handles payment
    refunds.
  \end{itemize}
\end{itemize}

\includegraphics[width=6.18593in,height=3.59524in]{media/image2.jpeg}

\subsection{Object Diagram}\label{object-diagram}

\subsection{Concept and Description}\label{concept-and-description}

The \textbf{Object Diagram} represents a snapshot of instances of
classes in a system at a particular moment. It helps visualize the
relationships between objects, showing their states and values.

\subsection{Application in Banking
System}\label{application-in-banking-system}

In a banking system, an Object Diagram can illustrate:

\begin{itemize}
\item
  An instance of a \texttt{Customer} object with linked \texttt{Account}
  objects.
\item
  An \texttt{Account} instance with active transactions.
\item
  A \texttt{Loan} instance tied to a specific \texttt{Customer}.
\end{itemize}

\subsection{Package Diagram}\label{package-diagram}

\subsection{Concept and Description}\label{concept-and-description-1}

A \textbf{Package Diagram} organizes the system's classes into logical
groups to manage complexity. It shows dependencies among different parts
of the system.

\subsection{Application in Banking
System}\label{application-in-banking-system-1}

\begin{itemize}
\item
  A package for \texttt{User\ Management} containing \texttt{Customer}
  and \texttt{BankEmployee} classes.
\item
  A package for \texttt{Transaction\ Processing} containing
  \texttt{Transaction}, \texttt{Payment}, and \texttt{Loan}.
\item
  A package for \texttt{Account\ Handling} containing \texttt{Account},
  \texttt{SavingsAccount}, and \texttt{CurrentAccount}.
\end{itemize}

\section{4. Behavioral Family
Diagrams}\label{behavioral-family-diagrams}

\subsection{4.1. Use Case Diagram}\label{use-case-diagram}

\subsubsection{Concept and Description}\label{concept-and-description-2}

The \textbf{Use Case Diagram} captures interactions between actors
(users) and the system. It highlights user goals and how the system
fulfills them.

\subsubsection{Application in Banking
System}\label{application-in-banking-system-2}

\begin{itemize}
\item
  \texttt{Customer} can open an account, make transactions, apply for
  loans, and check balance.
\item
  \texttt{BankEmployee} can approve loans, process transactions, and
  manage customers.
\end{itemize}

\subsection{4.2. Use Case Descriptions}\label{use-case-descriptions}

Each functionality has a description with pre-conditions and
post-conditions.

\begin{enumerate}
\def\labelenumi{\arabic{enumi}.}
\item
  \textbf{Open Account} \textbf{Pre-condition:} The user provides valid
  ID and initial deposit. \textbf{Post-condition:} A new account is
  created and linked to the customer.
\item
  \textbf{Deposit Money} \textbf{Pre-condition:} The customer has an
  active account. \textbf{Post-condition:} The account balance is
  updated.
\item
  \textbf{Withdraw Money} \textbf{Pre-condition:} Sufficient balance
  exists in the account. \textbf{Post-condition:} Money is deducted from
  the account.
\item
  \textbf{Apply for Loan} \textbf{Pre-condition:} The customer has a
  valid account. \textbf{Post-condition:} The loan application is stored
  for approval.
\item
  \textbf{Approve Loan} \textbf{Pre-condition:} The bank employee
  validates loan details. \textbf{Post-condition:} The loan status is
  updated to "Approved".
\end{enumerate}

\subsection{4.3 State Chart Diagram}\label{state-chart-diagram}

\subsubsection{Concept and Description}\label{concept-and-description-3}

A \textbf{State Chart Diagram} represents the various states an object
can be in during its lifecycle.

\subsubsection{Application in Banking
System}\label{application-in-banking-system-3}

\begin{itemize}
\item
  Account states: \texttt{Created} → \texttt{Active} → \texttt{Frozen} →
  \texttt{Closed}
\item
  Transaction states: \texttt{Initiated} → \texttt{Processing} →
  \texttt{Completed/Failed}
\item
  Loan states: \texttt{Applied} → \texttt{Approved} → \texttt{Disbursed}
  → \texttt{Repaid}
\end{itemize}

\subsection{4.4 Activity Diagram}\label{activity-diagram}

\subsubsection{Concept and Description}\label{concept-and-description-4}

An \textbf{Activity Diagram} represents workflows and processes in a
system.

\subsubsection{Application in Banking
System}\label{application-in-banking-system-4}

\begin{itemize}
\item
  The sequence of a fund transfer: \texttt{Start} →
  \texttt{Enter\ Details} → \texttt{Validate} →
  \texttt{Process\ Transaction} → \texttt{Update\ Balance} →
  \texttt{End}.
\end{itemize}

\subsection{4.5 Sequence Diagram}\label{sequence-diagram}

\subsubsection{Concept and Description}\label{concept-and-description-5}

A \textbf{Sequence Diagram} models interactions between objects in a
sequential order.

\subsubsection{Application in Banking
System}\label{application-in-banking-system-5}

\begin{itemize}
\item
  Customer initiates transaction → Bank verifies account → Transaction
  is processed → Confirmation is sent.
\end{itemize}

\subsection{4.6 Collaboration Diagram}\label{collaboration-diagram}

\subsubsection{Concept and Description}\label{concept-and-description-6}

A \textbf{Collaboration Diagram} represents interactions among objects
for a specific use case.

\subsubsection{Application in Banking
System}\label{application-in-banking-system-6}

\begin{itemize}
\item
  Customer interacts with ATM → ATM communicates with Bank System → Bank
  System updates account → ATM provides output.
\end{itemize}

Architectural Family Diagrams

\section{5. Architectural Family
Diagrams}\label{architectural-family-diagrams}

\subsection{5.1. Component Diagram}\label{component-diagram}

\subsubsection{Concept and Description}\label{concept-and-description-7}

A \textbf{Component Diagram} illustrates the physical structure of the
system, showing software components and their dependencies.

\subsubsection{Application in Banking
System}\label{application-in-banking-system-7}

\begin{itemize}
\item
  UI Layer - Web Interface, Mobile App.
\item
  Business Logic Layer - Account Management, Transaction Processing.
\item
  Data Layer - Database for storing customer, account, and transaction
  details.
\end{itemize}

\subsection{5.2. Deployment Diagram}\label{deployment-diagram}

\subsubsection{Concept and Description}\label{concept-and-description-8}

A \textbf{Deployment Diagram} represents hardware and software
deployment in the system.

\subsubsection{Application in Banking
System}\label{application-in-banking-system-8}

\begin{itemize}
\item
  Servers hosting the banking application.
\item
  Database server storing customer and transaction records.
\item
  ATM machines communicating with the server.
\end{itemize}

\section{Appendix}\label{appendix}

\subsection{Requirements Traceability
Matrix}\label{requirements-traceability-matrix}

\begin{longtable}[]{@{}
  >{\raggedright\arraybackslash}p{(\linewidth - 4\tabcolsep) * \real{0.2381}}
  >{\raggedright\arraybackslash}p{(\linewidth - 4\tabcolsep) * \real{0.4641}}
  >{\raggedright\arraybackslash}p{(\linewidth - 4\tabcolsep) * \real{0.2978}}@{}}
\toprule\noalign{}
\begin{minipage}[b]{\linewidth}\raggedright
Requirement ID
\end{minipage} & \begin{minipage}[b]{\linewidth}\raggedright
Description
\end{minipage} & \begin{minipage}[b]{\linewidth}\raggedright
Status
\end{minipage} \\
\midrule\noalign{}
\endhead
\bottomrule\noalign{}
\endlastfoot
RQ001 & Customer Account Management & Implemented \\
RQ002 & Online Fund Transfer & Implemented \\
RQ003 & Loan Approval Process & Under Development \\
RQ004 & Security Measures & Implemented \\
\end{longtable}

\subsection{Glossary of Terms}\label{glossary-of-terms}

\begin{itemize}
\item
  \textbf{Customer} - A user with a bank account.
\item
  \textbf{Transaction} - A financial operation (Deposit, Withdrawal,
  Transfer).
\item
  \textbf{Loan} - Borrowed money with an interest rate.
\item
  \textbf{ATM} - Automated Teller Machine for customer transactions.
\end{itemize}

\subsection{}\label{section}
